\chapter{A secret streamlet trickles (I First Met Her in a Dream... or Something)}
Puella Magi Madoka Magica can be understood as three distinct arcs, each focusing on a particular character or pair of characters and exploring a particular theme, although the themes of all three can be found in the other two. The first such arc covers the first three episodes, and focuses on Mami. Through her, it works through a fundamental internal struggle of the series, between its position within the magical girl genre and its aspirations to be something new. This in turn can be viewed as a struggle between two competing shows, the fairly typical (albeit somewhat dark), heavily Cardcaptor Sakura-influenced series it initially appears to be, and the deeply unsettling offspring of Revolutionary Girl Utena and Neon Genesis Evangelion it evolves into. To simplify still further, this can be understood as a conflict between a false show or mask (as epitomized by "Mata Ashita," the rather sweet and gentle ending credits sequence of the Blu-Ray version of the first two episodes) and a true show (as epitomized by the ending credits of Episode 3 on). In other words, what this first arc accomplishes is to set up a binary between the generic magical girl show audiences unfamiliar with series writer Gen Urobuchi likely expect, and the dark deconstructive series he is likely to create, before appearing to settle on the latter.

The first episode, and thus Madoka Magica itself, opens with a curtain rising. However, this is not traditional animation, but rather stop-motion animation of a paper curtain rising. Given that the difference between art and non-art is the frame-a story not framed as a story is a lie; a painting not framed as a painting is graffiti-a strong argument can be made that it is the frame that defines the art. This is, in essence, the institutionalist definition of art, that art is an artifact created by an artist and presented to an artworld public, that is, a public which has access to a theory of art and art history which permits them to recognize a work as art \cite{ref12}. In other words, art is that which codes itself as art in a way recognizable to its audience. In that case, we must consider the possibility that the stop-motion paper cutouts-and, more generally, the deployment of art styles far outside the norm for anime-are in some sense the definition of Madoka Magica, a representation of its individuality as opposed to its existence as an instance of a genre.

At the same time, the paper curtain evokes the Japanese tradition of kamishibai, paper theater, an early-20th century art form similar to Europe's traveling puppet shows, in which itinerant storytellers traveled from town to town, presenting tales accompanied by sequential illustrations displayed on a portable stage \cite{ref13}. This form, in turn, was a major influence on the development of manga and anime \cite{ref14}. Through the use of papercraft, Madoka Magica is simultaneously laying claim to the origins of its medium and genre, and asserting itself as an individual work independent of its genre.

This interpretation fits quite well, as we shall see, but first we must deal with Madoka running through a distorted, vast interior space, a checkerboard that evokes the warped spaces of Escher. Blacks and whites are sharply delineated here, clear binaries, but this space is not real. Unlike the latter part of Madoka's dream, it does not appear to correspond directly to events we see on any timeline, but rather is a hint at the place where she will shortly encounter her first witch.

The latter part of the dream, however, is (as we eventually learn near the end of the series) actually a memory of a previous timeline. This is the reality of the series, as signified by the use of "Magia," the series' true ending theme. "Magia" will not play at the ends of this or the next episode; in the TV broadcast, the credits play over the final scene of the episode, while the Blu-Ray version has a seemingly cheerful (so long as one does not pay much attention to the lyrics ) song accompanied by happy images. "Magia" is not played as ending credits until Episode 3, the point at which the series abandons all pretense of being a "normal" magical girl show.

Visually, this portion of the dream is all grays of various shades, with little black, white, or color. In the center of the ruined city is an enormous dead tree, nature and civilization fallen together. Walpurgisnacht appears, her Harlequin form and massive gears implying artificiality and order, belied by wild shrieks of laughter. Everything is its opposite. Duality is an illusion; in Madoka Magica all binaries collapse into unity. We will return here.

The episode tries to contain the revelations of this short sequence. The scene is played off as a dream, in much the same way as Cardcaptor Sakura-one of Madoka Magica's clearest influences-plays off its own portentous first scene \cite{ref15}. The opening credits which follow are also very typical for a magical girl show (again, so long as one ignores the lyrics), with Madoka's very Sailor Moon- \cite{ref16} or Cutie Honey-esque \cite{ref17} transformation sequence into a costume strongly reminiscent of several of Sakura's.

Indeed, most of the rest of the episode consists of playing up Madoka's life and personality to be as harmless and ordinary as possible. Her family is basically perfect, a warm and friendly-seeming stay-at-home gardener father and driven businesswoman mother, plus a typically adorable baby brother. She is nearly late for school and runs out with toast in mouth, just like the titular character in the first episode of Sailor Moon, among countless other examples. (Indeed, that precise image of a schoolgirl late for class, rushing out with toast, was used to convey the comfortable ordinariness of "normal" Rei in Shinji's fantasy sequence in the final episode of Neon Genesis Evangelion \cite{ref18}.)

Even the arrival of Homura, which shatters normality as far as Madoka herself is considered, is completely predictable from an audience perspective. The mysterious transfer student is a hoary enough anime trope to be invoked by name in the Haruhi Suzumiya novels \cite{ref19}, which are in part a parody of the several anime genres that share in common more or less typical Japanese adolescents encountering the supernatural; it thus borders on cliché that Homura appears following Madoka's dream, and quickly excels in every area of high school life. Her chilly persona (which, as a defensive affectation, fits both the classic and modern definitions of the word) and effortless, bored-seeming accomplishments serve only to make her more predictably mysterious. 

Save one, every intrusion of the world of the magical girls into Madoka's life is doubly generic-both not particularly distinctive, and typical of the genre. Her rescue of the injured Kyubey is yet another example, mirroring Sailor Moon's rescue of Luna \cite{ref20}, who likewise offered to unlock her potential in exchange for her fighting evil. Other than the one intrusion we're avoiding talking about, the most refreshingly non-generic moment in the episode is the intrusion of mundanity when Madoka describes her dream to Hitomi and Sayaka. In one of the most quietly brilliant scenes in the series, the three seriously discuss several explanations for the dream of various degrees of plausibility, from the reasonable suggestion that Madoka has seen Homura somewhere before to the absurdity of past life memory (which, of course, is the closest to the truth). Sayaka even notes how much like an anime the situation is in both the TV series and the film-a subtle foreshadowing that of all the characters to become magical girls, she is the one who is going to try hardest to uphold the idealized "warrior of justice" image that is the default anime depiction.

Which brings us to the least anime-like sequence in the episode, the witch attack. This is the ultimate intrusion of both the magical world into Madoka's safe little life, and of the non-anime into the anime. The stop-motion, papercraft explosion of bizarre and unsettling imagery forces Madoka and Sayaka to acknowledge the strangeness they had previously been shielded from, the wider, wilder world outside the one they know. The alien art style reinforces their disorientation and horror, while at the same time managing to be faintly absurd, even comical. The Anthonies in particular-mustachioed cotton balls reminiscent of the Pringles potato chip mascot-are both ridiculous and deeply unsettling. This is not an uncommon effect in postmodern works; the juxtaposition of elements from one context within another context creates a jarring, incongruous effect \cite{ref21}, and both humor \cite{ref22} and horror \cite{ref23} depend on similar sensations of incongruity.

Homura can do nothing about this intrusion except to say, "Oh no, not now." She recognizes that this is inevitable; as an agent of the true Madoka Magica (note that she is the only magical girl in Madoka's dream sequence) any attempt by her to preserve this safer, more comfortable false show is doomed to failure (as we will see again in both Episode 10 and Rebellion). That protection can only be accomplished-and only for a little while-by Homura's opposite number.

Enter Mami. 

She positions herself immediately as Homura's enemy, threatening her in order to protect Kyubey, who at this point in the story is still the cute, lovable mascot-herald who awakens the girls to their magic. It is only when the series drops its mask that he will drop his and become a manipulative devil-figure, taking over the antagonist role. Mami and Homura share in common, as we will see, that they are experienced magical girls who wield guns, but are otherwise near-total opposites. Homura is new to the school, while Mami, as an upperclassman, has been there longer than Madoka and Sayaka have. Homura's guns are entirely mundane, modern weapons, while Mami's are magical flintlocks. Homura is all straight lines, dark colors, and purples; Mami is all curves (not just in her figure, but her hair as well), whites, and yellows. Homura is closed, mysterious, seemingly hostile; Mami is open and friendly.

Most importantly, Mami has the power to restore the false Madoka Magica, where Homura does not. Her first act, before we even see her, is to define a safe space around Madoka and Sayaka within the witch's labyrinth. She is able to drive off the agents of the true Madoka Magica, both the witch and Homura. Much later in the series, we will see an alternate timeline where Mami learns the true nature of witches, one of the key differentiators between the "true" and "false" Madoka Magica, and her response is both violent and highly efficient, killing the ruthless magical girl Kyoko with her first shot, immobilizing Homura as her second action, and very nearly killing her before being killed herself by Madoka. She is ultimately put down, but at great cost, forcing that particular timeline (which can be viewed in a sense as a discarded "draft" version of the final series) to ultimately be abandoned-the "true" Madoka Magica simply does not work if she's around. 

Her role in this episode is to restore order, returning the art style to familiar anime norms, healing the injured Kyubey, and enabling him to take his initial position as the Luna-equivalent, offering magical power to Madoka and Sayaka. With her positive attitude, determination, and considerable power, Mami is a potent stand-in for magical girls past, and as we shall see over the course of this first arc, brings with her all the standard themes of the magical girl genre. So long as she stands in defense of it, the false series shall not fall.

She'll simply have to go.
